\section{Security model}

This project handles cryptographic key material and encrypted user data
directly. This section summarizes the trust boundaries described throughout
the manual so they are visible in one place.

\begin{securityBox}
Never commit private keys, tokens, secrets, API key bearers, provider
certificates, revocation lists, device descriptors (\file{device.kq},
\file{device.skey}), slot tokens (\file{*.kqst}), or plaintext copies of
protected files to a repository. See \file{.gitignore} for the patterns
already excluded. Compiling with \flag{--features provider} does not make a
host a trusted KeyQuorum provider.
\end{securityBox}

\subsection{What the mailbox can and cannot see}

The hosted mailbox cannot decrypt \file{.kqpb} envelopes and must not be
given wrapped shares or private keys. It stores the canonical \emph{public}
tree as JSON documents --- labels, fingerprints, public keys, policy.
Sending envelopes with \code{relay push} updates those documents from the
sender's \flag{--db}. Pulling returns a sliced copy for the pull-key
fingerprint, which the CLI translates into local SQLite. Personal devices
load a bearer with \code{keyquorum loadkey} (or \flag{--api-key} once);
they keep the hash and a sealed copy of the bearer in the owner-only org
database.

Recover a lost bearer by loading the replacement your provider issues.
Recover a missed update by pulling again and importing on the device that
holds the matching decryption key.

\subsection{Custody, quorum, and delegation, at a glance}

\begin{itemize}
  \item A key file with no container placement is its own device; distinct
  \flag{--share-file} keys are distinct physical devices, and presenting
  the same key twice still counts once (Section~\ref{sec:quorum-files},
  Section 3).
  \item Hardware custody requires each consumed Shamir share to sit on its
  own device id; logical custody lets several slots on one container meet a
  threshold together, but the container still counts as one device (Section 3).
  \item \code{unlock_approval = parent} layers a required parent signature
  on top of Shamir and the device count; it is opt-in, not a replacement for
  either (Section 3).
  \item A non-root \code{tree restructure} is a proposal until the parent
  countersigns; a direct tree update from that authorizer is refused; a
  routine employee reissue by the direct parent stays a single signature
  (Section~\ref{sec:private-bridges}, Section 9).
  \item A ghost identity keeps hierarchy and provenance but holds no usable
  private key: it cannot sign, satisfy a quorum, authorize an import, or be
  exported (Section~\ref{sec:device-transfer}).
\end{itemize}

\subsection{Short-lived caches and the relay check}

Three caches keep you from repeating work (Section~\ref{sec:caches}):
the parameters you just used, a relay identity check that just passed, and
facts \code{doctor} already looked up. They share one flat 15-minute
lifetime, hold no passphrase, key, bearer or plaintext, and are never an input
to a signature, a quorum, custody, an approval, a freshness judgment or a
trust decision; each of those is computed again from the store every time.

The relay entry is the one that stands in for a challenge. A relay proves a
KeyQuorum-signed certificate and signs a random challenge; a cached pass
instead reuses that earlier result for the same URL while the revocation list
and the stored key hash are unchanged, for at most 15 minutes and never past
the expiry of the certificate that was checked. It is never
written for a failure, and \code{loadkey} always runs the full challenge.
What it does not do is prove who answers \emph{now}: for up to 15 minutes a
reused pass trusts the earlier proof of that URL, so someone who took over the
endpoint in that window would be sent the stored bearer. Letters stay sealed
and signed either way. Where that window matters, use \flag{--no-cache},
\code{KEYQUORUM_NO_CACHE=1} or \code{use --cache off}, and every command
challenges the relay again.

\subsection{Envelopes and authenticity}

Every sealed letter in this manual --- a private-bridge package, an
authenticated update, an export bundle, a delivered file, a device transfer
letter --- shares the same shape: the outside names a recipient public key
in the clear, the inside is sealed so only that recipient's private key can
open it, and where the letter carries an authorization decision (an
authenticated update, a delivery, a device transfer acknowledgement) it is
additionally \emph{signed} by whoever is authorizing that decision, checked
against a key the opening store already has on file. A mailbox that carries
these letters can route them by their outer fingerprint without ever being
able to read or forge what is inside.

\subsection{Secrets in memory and the relay's audit trail}

Passphrases, passwords, PINs and API keys typed at a prompt, the stored relay
bearer, and every decrypted plaintext or unsealed letter are held in buffers
that are overwritten when the command is done with them. What a command writes
out on purpose (to stdout, or a file you name) is yours to protect.

The relay records each API key's creation, rotation and revocation, and every
attempt to mint one, in a hash chain whose head the relay signs with its
KeyQuorum-certified key. That signature is trusted only for the period the
relay's certificate was valid, and not at all once the certificate is revoked.
A key holder can read the events about their own key, and nobody else's.

